% !PRINTED 116
\weberpage{116}
If one further sets $u = 0$ in (7), (8), (9), \S\,33, then uses the relations (10) and (20) and extracts the square root, one finds

\weq{21}{\begin{aligned}
f(n\,\omega)\,\eta(\omega)^{\frac{n-1}{2}} &= f(\omega) \prod_{1,\,\frac{n-1}{2}}^{\nu} \vartheta_{00}\!\left(\frac{\nu}{n}\right), \\
f_1(n\,\omega)\,\eta(\omega)^{\frac{n-1}{2}} &= f_1(\omega) \prod_{1,\,\frac{n-1}{2}}^{\nu} \vartheta_{01}\!\left(\frac{\nu}{n}\right), \\
f_2(n\,\omega)\,\eta(\omega)^{\frac{n-1}{2}} &= f_2(\omega) \prod_{1,\,\frac{n-1}{2}}^{\nu} \vartheta_{10}\!\left(\frac{\nu}{n}\right).
\end{aligned}}

The signs follow here from the assumption of an infinitely small $q$.

\begin{center}
\textbf{\S\,35. \gesperrt{The Weierstrass $\sigma$-Function.}}
\end{center}

By the general linear substitution
\[
\begin{pmatrix} \alpha, & \beta \\ \gamma, & \delta \end{pmatrix}, \qquad \alpha\delta - \beta\gamma = 1,
\]
applied to the variables $\omega_1$, $\omega_2$, every $t$-function passes again into a $t$-function, whereby the characteristic has changed. If $(g_1, g_2)$ is the characteristic of the original function and $(g_1', g_2')$ that of the transformed function, then by \S\,27, (11)
\[
(g_1, g_2) = (\delta g_1' - \beta g_2' - \beta\delta,\ -\gamma g_1' + \alpha g_2' - \alpha\gamma).
\]

If one solves the two linear equations contained in this for $g_1'$, $g_2'$, and observes that $\alpha\beta(\gamma+\delta+1)$, $\gamma\delta(\alpha+\beta+1)$ are necessarily even numbers, and that a characteristic does not change when its elements change by multiples of 2, then one may also write in place of it:

\weq{1}{(g_1', g_2') = (\alpha g_1 + \beta g_2 + \alpha\beta,\ \gamma g_1 + \delta g_2 + \gamma\delta).}

The $\vartheta$-functions therefore pass into one another, apart from exponential factors. It is however of importance to form a function which is \gesperrt{absolutely invariant} with respect to the linear transformations, and such a function is the $\sigma$-function introduced into the theory by \gesperrt{Weierstrass}, to whose definition we now proceed.

If we apply formula (1) to the four principal characteristics: $(0, 0)$, $(0, 1)$, $(1, 0)$, $(1, 1)$, then these pass over, in order,
%CONT>>
